\subsection{RL Teacher Training and Safety Details}
\label{app:rl_training_details}

Table~\ref{tab:teacher_observations} contrasts the actor's noisy, deployable
signals with the centralized critic's clean, privileged simulator state.
The palm and five fingertips constitute the six hand links. A ``twist'' is the
concatenation of three-dimensional linear and angular velocity. Running
mean/variance normalization is applied to both inputs and to value targets.

\begin{table}[H]
  \centering
  \caption{Asymmetric teacher observations.}
  \label{tab:teacher_observations}
  \scriptsize
  \begin{minipage}{0.94\textwidth}
  \setlength{\tabcolsep}{2.8pt}
  \renewcommand{\arraystretch}{1.04}
  \begin{tabularx}{\linewidth}{@{}
    p{0.25\linewidth}
    >{\centering\arraybackslash}p{0.10\linewidth}
    >{\centering\arraybackslash}p{0.07\linewidth}
    X@{}}
    \toprule
    Signal & Actor & Critic & Definition \\
    \midrule
    Robot joint position $q_t$ & 13 (noisy) & 13 & Actuated arm--hand joint angles (rad) \\
    Robot joint velocity $\dot q_t$ & 13 (noisy) & 13 & Actuated arm--hand joint velocities (rad/s) \\
    Hand-link position & 18 (noisy) & 18 & Palm and five fingertip positions: $6\times3$ Cartesian coordinates \\
    Hand-link velocity & 18 (noisy) & 36 & Actor: linear velocity ($6\times3$); critic: spatial twist ($6\times6$) \\
    Palm contact force & -- & 3 & Privileged Cartesian contact-force vector (N) \\
    Measured joint torque & -- & 19 & Privileged torque measurements for 19 instrumented joints (N\,m) \\
    Object position & 3 (noisy) & 3 & Object Cartesian position \\
    Object orientation & 4 (noisy) & 4 & Object rotation as a unit quaternion \\
    Object velocity & -- & 6 & Privileged object linear and angular velocity (twist) \\
    Goal position & 3 & 3 & Cartesian target position for the lift task \\
    Object scale & 1 & 1 & Scalar geometry scale supplied by the environment \\
    Object ID & 1 & 1 & Scalar identifier for the current object \\
    Previous action $a_{t-1}$ & 13 & 13 & Previous 13-dimensional arm--hand command \\
    \midrule
    \textbf{Total dimension} & \textbf{87} & \textbf{133} & \\
    \bottomrule
  \end{tabularx}

  \vspace{2pt}
  \raggedright\scriptsize
  ``--'' denotes a signal that is unavailable to the actor. Noisy entries are
  sensor-like estimates; unmarked actor entries are task context or action
  history. Critic-only entries are used during simulation training and are not
  available to the deployed policy.
\end{minipage}
\end{table}

\paragraph{Termination conditions.}
\placeholder{Report the numerical trigger, units, measured signal, and physical
justification for harmful collision, object escape, palm flipping, and the
hand-too-far/workspace condition. State how simultaneous triggers are resolved
and whether the same definitions are used for teacher training, distillation,
and evaluation. Report the available contact-force or other severity signal and
define the severity-weighted safety metric.}
\requestedby{Reviewer \#22}

\paragraph{Runs and random seeds.}
Teacher performance in Main Paper Table~II is evaluated over 200 episodes per
object, and the real-world results are evaluated over 30 episodes per object.
Each of the 40 object-specific teacher policies is trained with five independent
random seeds $\{1,2,3,4,5\}$, with the same evaluation protocol applied to every
seed. For each seed, teacher performance is evaluated on all 40 objects. The
table below reports the mean~$\pm$~standard deviation across the 40 object-level
rates for lift success, overall unsafe events, and the four unsafe-event
categories separately for each seed.
\requestedby{Reviewers \#14 and \#22}

\begin{table}[H]
  \centering
  \small
  \caption{Teacher policy performance across all 40 objects for each random seed
  (200 evaluation episodes per object and seed). Each entry is the mean
  $\pm$ standard deviation across the 40 object-level rates. Numerical entries
  are provisional synthetic placeholders pending verification with the measured
  evaluation data.}
  \label{tab:seed_reporting}
  \setlength{\tabcolsep}{9pt}
  \renewcommand{\arraystretch}{1.14}
  \begin{tabular}{l|ccccc}
    \toprule
    Metric (\%) & Seed 1 & Seed 2 & Seed 3 & Seed 4 & Seed 5 \\
    \midrule
    Lift success      & 86.7 $\pm$ 7.8 & 87.4 $\pm$ 7.5 & 86.9 $\pm$ 8.0 & 87.8 $\pm$ 7.3 & 87.1 $\pm$ 7.7 \\
    Unsafe rate avg.  & 14.2 $\pm$ 6.1 & 13.6 $\pm$ 5.8 & 14.0 $\pm$ 6.3 & 13.2 $\pm$ 5.6 & 13.8 $\pm$ 6.0 \\
    Harmful collision & 3.1 $\pm$ 3.5 & 2.9 $\pm$ 3.3 & 3.2 $\pm$ 3.6 & 2.8 $\pm$ 3.2 & 3.0 $\pm$ 3.4 \\
    Object escape     & 9.0 $\pm$ 5.2 & 8.8 $\pm$ 5.0 & 8.9 $\pm$ 5.3 & 8.7 $\pm$ 4.9 & 8.9 $\pm$ 5.1 \\
    Palm flipped      & 1.2 $\pm$ 2.4 & 1.1 $\pm$ 2.3 & 1.1 $\pm$ 2.3 & 1.0 $\pm$ 2.2 & 1.1 $\pm$ 2.3 \\
    Hand too far      & 0.9 $\pm$ 1.9 & 0.8 $\pm$ 1.8 & 0.8 $\pm$ 1.8 & 0.7 $\pm$ 1.7 & 0.8 $\pm$ 1.8 \\
    \bottomrule
  \end{tabular}
\end{table}

\paragraph{Teacher safety and disagreement handling.}
Main Paper Table~II reports representative teacher unsafe-event rates. Across
all 40 teacher policies, unsafe-event rates range from 5--30\%. Because an
intervention transfers control to an imperfect teacher, the framework's safety
performance is restricted by the teacher's own performance ceiling. Teacher
takeover can reduce unsafe exploration, but it does not provide a formal safety
guarantee. \placeholder{Define how the controller handles disagreement between
the teacher action and the critic's assessment, including the fallback used
when neither candidate action satisfies the acceptance rule.}
\requestedby{Reviewer \#22}

\subsection{VLM Prompt and Inference Details}
\label{app:vlm_inference_details}

The VLM prompt specifies the system and user requirements, intervention rule,
multimodal input format, task description, and output keys, as shown below.

\begin{algorithm}[H]
\makeatletter
\renewcommand{\fnum@algorithm}{}
\makeatother
\caption{\textbf{VLM Prompt Example}}
{\small\ttfamily\renewcommand{\baselinestretch}{1.12}\selectfont
\begin{algorithmic}
\setlength{\itemsep}{1pt}
\STATE \textbf{System:}
\STATE You are a robotics safety assistant. Your task is to assess the current
safety--autonomy operating point of a visuomotor policy.
\STATE \textbf{Inputs:}
\STATE \hspace{1em}- Current visual observation: \textless IMAGE\textgreater
\STATE \hspace{1em}- Teacher--student disagreement: \{d\}
\STATE \hspace{1em}- Predicted success probability: \{q\_succ\}
\STATE \hspace{1em}- Intervention rate: \{r\_int\}
\STATE \hspace{1em}- Unsafe-event rate: \{r\_unsafe\}
\STATE \textbf{Task:}
\STATE Determine whether the current intervention boundary should become more
conservative or more permissive.
\STATE \textbf{Return:}
\STATE \{
\STATE \hspace{1em}"disagreement\_threshold": \ldots,
\STATE \hspace{1em}"success\_threshold": \ldots,
\STATE \hspace{1em}"reason": \ldots\ \}
\end{algorithmic}
}
\end{algorithm}

\requestedby{Reviewers \#14 and \#22}

\paragraph{Model and input configuration.}
We used the OpenAI GPT-5.5 multimodal model through the Responses API
(\texttt{model="gpt-5.5"}) with reasoning effort set to \texttt{high}
(\texttt{reasoning.effort="high"}). The model received both textual rollout
statistics and RGB images. All other generation parameters were left at their
API defaults. We issued the first VLM query after 2{,}000 episodes and
subsequently queried it every 10{,}000 distillation episodes, resulting in 10
VLM calls over the full 100{,}000-episode training run and approximately
70{,}000 tokens of total usage per run.
\requestedby{Reviewers \#14 and \#22}

\paragraph{Output processing and threshold updates.}
\placeholder{Provide the machine-readable output schema, parsing and validation
logic, permitted threshold ranges, clipping and smoothing rules, and query/update
frequency. Explain exactly how the parsed output changes the intervention
policy. State explicitly whether VLM adaptation is used only during simulation
distillation, during real-world student training, or during physical deployment.}
\requestedby{Reviewers \#14 and \#22}

\paragraph{Reliability and resource usage.}
\placeholder{Report retry and timeout behavior, malformed-output handling,
deterministic fallback, measured latency, token usage, number of calls, and
monetary cost per training run.}
\requestedby{Reviewer \#22}

\clearpage
\subsection{Additional Experimental Results}
\label{app:additional_experimental_results}

\paragraph{Bellman critic evaluation.}
\placeholder{State how safety risk and task success are separated, or justify
the retained target. Report AUROC, AUPRC, Brier score, expected calibration
error, reliability diagrams, and the false-negative rate for unsafe episodes at
the deployed threshold.}
\requestedby{Reviewers \#14 and \#22}

\paragraph{VLM-necessity and threshold-adaptation ablations.}
\placeholder{Report matched comparisons for fixed thresholds, grid-search and
Bayesian-optimization selection, rule-based adaptation, curriculum scheduling,
a small learned threshold controller, an LLM receiving rollout statistics only,
the VLM receiving images only, the VLM receiving statistics and images, and a
visual-input removal or mismatch control. Include task success, unsafe-event
rate, and intervention cost under matched budgets.}
\requestedby{Reviewers \#14 and \#22}

\paragraph{Training and resource comparison.}
\placeholder{For Vanilla DAgger, SafeDAgger, critic-based intervention, and
VLM-guided intervention, report environment steps, student updates, teacher
queries or labels, teacher-controlled steps, intervention rate, compute time,
VLM calls and cost, final success and unsafe-event rates, and results over
multiple seeds.}
\requestedby{Reviewer \#22}

\paragraph{Complete real-world evaluation.}
We will report results for the remaining 16 real-world test objects that are not
shown in the manuscript's six-object table, using the same success and safety
breakdown. The two eight-object blocks below separate task success from the
observed unsafe failure modes. Object names and numerical results are retained
as placeholders pending the expanded evaluation.
\requestedby{Reviewers \#14 and \#22}

\begin{table}[H]
\centering
\scriptsize
\caption{Template for per-object experimental results for the remaining 16
objects. Object names and numerical entries are placeholders.}
\label{tab:additional_experimental_results}
\setlength{\tabcolsep}{3pt}
\renewcommand{\arraystretch}{1.08}
\resizebox{0.90\textwidth}{!}{%
\begin{tabular}{l|cccccccc}
\hline
Metric (\%) & Object 7 & Object 8 & Object 9 & Object 10 & Object 11 & Object 12 & Object 13 & Object 14 \\
\hline
Lift Success      & -- & -- & -- & -- & -- & -- & -- & -- \\
Unsafe Rate Avg   & -- & -- & -- & -- & -- & -- & -- & -- \\
Harmful collision & -- & -- & -- & -- & -- & -- & -- & -- \\
Object escape     & -- & -- & -- & -- & -- & -- & -- & -- \\
Palm flipped      & -- & -- & -- & -- & -- & -- & -- & -- \\
Hand too far      & -- & -- & -- & -- & -- & -- & -- & -- \\
\hline
\end{tabular}%
}

\vspace{0.6em}
\resizebox{0.90\textwidth}{!}{%
\begin{tabular}{l|cccccccc}
\hline
Metric (\%) & Object 15 & Object 16 & Object 17 & Object 18 & Object 19 & Object 20 & Object 21 & Object 22 \\
\hline
Lift Success      & -- & -- & -- & -- & -- & -- & -- & -- \\
Unsafe Rate Avg   & -- & -- & -- & -- & -- & -- & -- & -- \\
Harmful collision & -- & -- & -- & -- & -- & -- & -- & -- \\
Object escape     & -- & -- & -- & -- & -- & -- & -- & -- \\
Palm flipped      & -- & -- & -- & -- & -- & -- & -- & -- \\
Hand too far      & -- & -- & -- & -- & -- & -- & -- & -- \\
\hline
\end{tabular}%
}
\end{table}

\placeholder{Replace the object labels and dashes with the remaining 16
per-object results. State 30 episodes per object, report 95\% confidence
intervals and statistical comparisons, and provide aggregate results separately
for seen and unseen objects across the full 22-object evaluation.}
\requestedby{Reviewers \#14 and \#22}

\paragraph{Physical baseline comparison.}
\placeholder{Report the feasible matched real-robot baseline using the same
objects, initial-condition protocol, and number of trials, including lift
success, unsafe-event incidence, uncertainty, and the statistical comparison to
VLM-DexSafeDagger.}
\requestedby{Reviewer \#14}
